\textbf{Ionizace} je vznik iontů; odtržením elektronu z neutrálního atomu vznikne kladný iont a volný elektron. Je třeba dodat alespoň ionizační energii $E_i$.
\begin{itemize}
\item \textbf{Nárazová ionizace:} srážka s částicí, jejíž kinetická energie splňuje $\tfrac12mv^2\ge E_i$.
\item \textbf{Fotoionizace:} působení záření s energií fotonu $hf\ge E_i$; zdrojem může být UV, rentgenové či gama záření.
\item \textbf{Tepelná ionizace:} dostatečně vysoká teplota zajistí energii pro ionizační srážky.
\end{itemize}
Podle působení rozlišujeme \textbf{objemové zdroje} (záření, teplo v objemu plynu) a \textbf{povrchové zdroje} (uvolňování elektronů z elektrod).
